\cgeentry{cge_macro}{Macros of the language of user variables and equations}

\cgesec{Description}

\cmd{cge_macro} (from version 0.9.7) defines a macro of the language of
\manref{cge_variable} and \manref{cge_equation}, as GAMS's \cmd{\$macro} does:
a name with parameters that stands for an expression. A CES aggregate, a
price deflated by the consumer price index or a constant written once can then
be used wherever an expression may be, with the arguments in place of the
parameters. \cmd{cge_macro, list} displays the macros, and
\cmd{cge_macro, clear} removes them.

\cgesec{Quick start}

A CES aggregate of two quantities
\begin{stcode}
cge_macro ces(x, y, d, r) = (d * x^(-r) + (1 - d) * y^(-r))^(-1 / r)
\end{stcode}
Use it in a user variable over the commodities that are imported
\begin{stcode}
cge_variable CQ(c)$(base(QM(c)) > 0) = ces(QD(c), QM(c), 0.5, 0.5)
\end{stcode}
A macro without parameters, and one that uses others
\begin{stcode}
cge_macro half = 0.5
cge_macro cesq(c) = ces(QD(c), QM(c), half, half)
\end{stcode}
Display the macros; remove them
\begin{stcode}
cge_macro, list
cge_macro, clear
\end{stcode}

\cgesec{Syntax}

Define a macro
\begin{stsyntax}
\cmdb{cge_macro} \argu{name}[\cmd{(}\argu{a}[\cmd{,} \argu{b} \dots]\cmd{)}] \cmd{=} \argu{expression}
\end{stsyntax}
Display or remove the macros
\begin{stsyntax}
\cmdb{cge_macro} \texttt{,} \opt{}{list}\\
\cmdb{cge_macro} \texttt{,} \opt{}{clear}
\end{stsyntax}
\argu{name} is a new name: not a variable, a parameter, an element, a set, a
user equation, a function of the language or another macro. The parameters
\argu{a}, \argu{b}, \dots\ are names, each given once. \argu{expression} is an
expression of the language of \manref{cge_equation}, in which the parameters
stand for the arguments.

\begin{opttable}
\optgroup{Main}
\opt{}{list} & display the macros\\
\opt{}{clear} & remove the macros\\
\end{opttable}

\cgesec{Options}

\begin{optdesc}
\item[\opt{}{list}] displays the macros with their parameters and expressions.

\item[\opt{}{clear}] removes every macro. A macro that a user variable or
equation uses cannot be removed until they are.
\end{optdesc}

\cgesec{Remarks and examples}

In an expression, \argu{name}\cmd{(}\argu{x}\cmd{,} \argu{y}\cmd{)} is replaced by
the macro's expression, in parentheses, with \argu{x} in place of \argu{a}
and \argu{y} in place of \argu{b}. An argument that is more than a name or a
number is itself put in parentheses, so that with \cmd{sq(x) = x * x},
\cmd{sq(QA(a) + 1)} is \cmd{((QA(a) + 1) * (QA(a) + 1))}, not
\cmd{QA(a) + 1 * QA(a) + 1}. A macro without parameters stands for its
expression. The expression may use variables, parameters, user sets and
parameters, sums, conditions, the functions of the language and other macros,
which are replaced in turn; a macro may not use itself, directly or through
others. Arguments may be indexes: in the definition of a variable over
\cmd{c}, \cmd{cesq(c)} is written out for every commodity.

Macros are replaced when a variable or equation is added, before its indexes
are written out, and again each time the model is calibrated. They are saved
with the model (\manref{cge_save}) and kept when it is calibrated again. In
the GAMS program, \manref{cge_gams} writes the expressions out, so it needs no
\cmd{\$macro}.

The SYNTHETIC SAM is mapped and calibrated:
\stlog{mac_setup}

\cgesubsec{Example 1: macros and their use}

A CES aggregate, a price relative to the consumer price index, a constant,
and a macro that uses the others; then variables defined with them, after the
tariffs are removed:
\stlog{mac_define}
\cmd{CQ} is $\big(0.5\,\V{QD}^{-0.5} + 0.5\,\V{QM}^{-0.5}\big)^{-2}$ for each
commodity with imports at the base, and \cmd{RP} the composite price
deflated by the consumer price index.

\cgesubsec{Example 2: errors}

A macro needs its arguments, as many as its parameters; a macro may not use
itself; and a macro that variables use cannot be removed:
\stlog{mac_err}

\cgesec{Stored results}

\cmd{cge_macro} stores the following in \cmd{r()}:
\begin{storedres}
\resgroup{Scalars}
\res{r(N_macros)} & number of macros\\
\resgroup{Macros}
\res{r(name)} & name of the macro defined\\
\end{storedres}

\cgesec{Methods and formulas}

The tokens of an expression are scanned for the names of macros. Each use,
\argu{name}\cmd{(}$x_1$\cmd{,} \dots\cmd{,} $x_m$\cmd{)}, is replaced by the
tokens of the macro's expression, with each parameter replaced by its
argument, in parentheses unless the argument is a single token, and the
whole in parentheses. The scan is repeated while macros remain, up to 30
times; a chain of macros that does not end within that is an error. The
macros are replaced before \cmd{card()} and before the indexes and sums are
written out (\manref{cge_equation}), so that the indexes in a macro's
arguments are written out with the others. A macro therefore adds no
variables, equations or derivatives of its own.

\cgesec{Also see}

\begin{alsosee}
\seeentry{cge_equation}{User equations: new relations among the model's variables}
\seeentry{cge_variable}{User variables, defined by an expression or determined by user equations}
\seeentry{cge_set}{User sets of elements, for indexes and sums}
\seechapter{chap:user}{User variables and equations}
\end{alsosee}
